\documentclass[10pt,a4paper]{amsart}
\usepackage[margin=19mm]{geometry}
\usepackage{amsmath,amssymb,mathtools}
\usepackage[scr=rsfs,cal=euler]{mathalfa}
\usepackage{xfrac}
\usepackage[unicode,hidelinks,bookmarks=false]{hyperref}
\newcommand{\ie}{i.e.\ }
\newcommand{\cf}{cf.\ }
\newcommand{\resp}{respectively\ }
\newcommand{\Subsection}{\subsection}
\providecommand{\nobreakdash}{\nobreak\hbox{-}}
\setlength{\emergencystretch}{2em}
\multlinegap0pt
\usepackage{bm}
\usepackage{bbm}
\usepackage{dsfont}
\usepackage{xspace}
\usepackage{cancel}
\usepackage{mathtools}
\usepackage{amsthm}
\makeatletter
\@ifundefined{theorem}{\newtheorem{theorem}{Theorem}}{}
\@ifundefined{lemma}{\newtheorem{lemma}[theorem]{Lemma}}{}
\makeatother
\def\B{{B}}
\def\N{{\mathbb N}}
\def\S{{\mathbb S}}
\def\dist{{\rm dist\,}}
\def\eps{\varepsilon}
\def\n{{\mathcal M}}
\newcommand{\R}{\mathbb{R}}
\newcommand{\abs}[1]{\left |#1 \right |}
\newcommand{\ageq}{\succsim}
\newcommand{\aleq}{\precsim}
\newcommand{\barint}{
\rule[.036in]{.12in}{.009in}\kern-.16in \displaystyle\int }
\newcommand{\brac}[1]{\left (#1 \right )}
\newcommand{\dif}{\,\mathrm{d}}
\newcommand{\dx}{\dif x}
\title[Existence of nontrival $n$-harmonic maps via min-max methods]{Existence of nontrival $n$-harmonic maps via min-max methods}
\author{Dorian Martino, Katarzyna Mazowiecka, Armin Schikorra}
\date{Source: arXiv:2601.05700v1 (CC BY 4.0)}
\begin{document}
\maketitle

\noindent\emph{Source and licence.} Existence of nontrival $n$-harmonic maps via min-max methods. Version arXiv:2601.05700v1 (CC BY 4.0), distributed under the Creative Commons Attribution 4.0 International licence, as recorded in the arXiv licence metadata for this exact version and shown on the abstract page https://arxiv.org/abs/2601.05700v1. Full rights notice: licenses/arxiv-2601.05700-CC-BY-4.0.txt.\par\smallskip

\noindent\textit{Editorial scope.} Counted unit: the Lemma whose source label is \texttt{la:minimalbeta} in the article (source0.tex lines 411--414, source label \texttt{la:minimalbeta}), delivered together with the article's own complete original proof of it (source0.tex lines 415--487, one \texttt{proof} environment of 73 lines). No mathematics is written, repaired or re-derived: statement and proof are the article's own text line for line. Required context retained uncounted: eq:betapvvalue (lines 407--409). Every definition, display and label of those lines is retained unchanged. Omitted from this package and not redistributed: 1--406, 410--410, 488--951. Nothing inside a retained excerpt is omitted or altered: every retained body is delivered exactly as the article's lines are written, so no omission receipt of any kind accompanies this package. Citations: the retained excerpts carry no citation command at all, so no bibliography entry of the article is transcribed and no reference is redistributed. Reference adaptation: none was needed and none was made; every reference inside the delivered unit resolves inside it, so no unresolved reference or citation is produced. Printed slips kept literal: no printed word, symbol, hypothesis or number of the counted statement or of its proof was altered. Layout and class: the article's own class is \texttt{amsart} with packages including pgf, tikz, pgfplots, mathrsfs, hyperref, backref, cleveref, xcolor, amsthm, latexsym, amsmath, amssymb, accents, todonotes; the wrapper is the lane's pinned \texttt{amsart} editorial wrapper at 10pt on A4, which restates the theorem-like environments the retained text needs and adds the article's own macro definitions that the retained text needs (\texttt{\textbackslash B}, \texttt{\textbackslash N}, \texttt{\textbackslash R}, \texttt{\textbackslash S}, \texttt{\textbackslash abs}, \texttt{\textbackslash ageq}, \texttt{\textbackslash aleq}, \texttt{\textbackslash barint}, \texttt{\textbackslash brac}, \texttt{\textbackslash dif}, \texttt{\textbackslash dist}, \texttt{\textbackslash dx}, \texttt{\textbackslash eps}, \texttt{\textbackslash n}), with extra wrapper packages (bm, bbm, dsfont, xspace, cancel, mathtools). The delivered numbering is the unit's own. Assets: no retained span contains a figure environment, a graphics-inclusion command, a PDF-page inclusion, a table or a TikZ picture; no image file is copied into this package, the package declares no asset dependency and no image file is redistributed with it. Licence: version arXiv:2601.05700v1 of this article is distributed under the Creative Commons Attribution 4.0 International licence, as recorded in the arXiv licence metadata for this exact version and shown on the abstract page https://arxiv.org/abs/2601.05700v1. A full rights notice is delivered at licenses/arxiv-2601.05700-CC-BY-4.0.txt. Characters: every retained excerpt is pure ASCII, so the delivered engine, XeLaTeX with its default Latin Modern text font, reports no missing character.
\begin{equation}\label{eq:betapvvalue}
 \inf_{\Phi \in \alpha}\ \max_{t \in [-1,1]^k}\ \mathcal{E}_p(\Phi(\cdot,t))=\inf_{\Phi \in \alpha}\ \max_{t \in [-1,1]^k}\ \int_{\S^n} |\nabla \Phi(x,t)|^p  \dx\eqqcolon \beta_p.
\end{equation}

\begin{lemma}\label{la:minimalbeta}
Let $\beta_p$ be the value defined in \eqref{eq:betapvvalue}. We have
\[\liminf_{p \to n^+} \beta_p > 0. \]
\end{lemma}

\begin{proof}
By H\"older inequality it suffices to show
\[
 |\S^n|^{1-\frac{p}{n}} (\beta_p)^{\frac{n}{p}} \geq \inf_{\Phi \in \alpha} \max_{t \in [-1,1]^k} \int_{\S^n} |\nabla \Phi(x,t)|^n \dx > 0.
\]
Assume on the contrary that for a suitably small number $\eps > 0$ (chosen below) there exists a smooth $\Phi \colon \S^{n}\times[-1,1]^k \to \n$ for which we have
\begin{equation}\label{eq:smallness-contradiction}
 \max_{t \in [-1,1]^k} \int_{\S^n} |\nabla \Phi(x,t)|^n \dx < \eps.
\end{equation}
For each $t \in [-1,1]^k$, let $\Phi^h(\cdot,t)$ denote the harmonic extension of $\Phi(\cdot,t)\colon \S^n \to \n$  to a map $\Phi^h(\cdot,t)\colon \B^{n+1} \to \R^N$.
If $P(x,\xi) = \frac{1-|x|^2}{\omega_n |x-\xi|^{n+1}}$ is the Poisson kernel on the ball $B^{n+1}$, then we have 
\[
 \Phi^h(x,t) = \int_{\S^n} P(x,\xi) \Phi(\xi,t)\dif \xi.
\]
Since we have $\int_{\S^n}P(x,\xi)\dif \xi = 1$, the following holds for any $\zeta \in \S^n$
\[
\begin{split}
 \dist(\Phi^h(x,t),\n) &\le \abs{\Phi^h(x,t) - \Phi(\zeta,t)} \le \int_{\S^n} P(x,\xi) \abs{\Phi(\xi,t) - \Phi(\zeta,t)}\dif \xi.
\end{split}
 \]
Multiplying both sides by $P(x,\zeta)$ and integrating over $\S^n$ with respect to $\zeta$ we obtain
\begin{equation}\label{eq:distanceestimate}
 \begin{split}
  \dist(\Phi^h(x,t),\n) &\le \int_{\S^n}\int_{\S^n} P(x,\xi) P(x,\zeta) \abs{\Phi(\xi,t) - \Phi(\zeta,t)}\dif \xi \dif \zeta. 
 \end{split}
\end{equation}
If $|x|\le \frac12$, then we have $|x-\xi|\ge \frac12$ for any $\xi\in \S^n$ and thus, it holds
\[
 \int_{\S^n}|P(x,\xi)|^{n'}\dif \xi \aleq \int_{\S^n}\brac{\frac{1}{|x-\xi|^{n+1}}}^{n'} \dif \xi \aleq 1.
\]
Hence, for $|x|\le \frac12$, we have, using H\"{o}lder and Poincar\'{e} inequality
\[
\begin{split}
 \dist(\Phi^h(x,t),\n) 
 &\le \|P(x,\cdot)\|_{L^{n'}(\S^n)}^2 \brac{\int_{\S^n}\int_{\S^n} \abs{\Phi(\xi,t) - \Phi(\zeta,t)}^n\dif \xi \dif \zeta}^{\frac 1n}\\
 &\aleq \brac{\int_{\S^n}|\nabla \Phi(x,t)|^n \dif \xi}^{\frac 1n} \aleq \eps^{\frac 1n}.
\end{split}
 \]
If $|x|\ge \frac12$, then we let $r=1-|x|$ and $X=\frac{x}{|x|}$ and we write 
\[
A(X,r,0) \coloneqq B(X,r), \quad  A(X,r,k)\coloneqq B(X,2^{k}r)\setminus B(X,2^{k-1}r) \text{ for }k\ge 1.
\]
Let us denote by $\theta = \dist_{\S^n}(\frac{x}{|x|},\xi)$. Then, we have
\[
 |x-\xi|^2 = |x|^2+1-2|x|\frac{x}{|x|}\cdot \xi = |x|^2+1-2|x|\cos\theta = (1-|x|)^2+2|x|(1-\cos\theta).
 \]
For every $\theta\in[0,1]$ we have the elementary inequality $\frac{2}{\pi^2}\theta^2\leq 1-\cos \theta$, thus
\[
 |x-\xi|^2\ageq (1-|x|)^2+|x|\theta^2 \ageq (1-|x|)^2 + \theta^2.
\]
Hence, for $\xi \in A(X,r,k)$ we have $\theta \approx 2^k r$ and $|x-\xi|\ageq 2^k r$. Thus, from \eqref{eq:distanceestimate} we get for $|x|\ge \frac 12$
\[
\begin{split}
 \dist(\Phi^h(x,t),\n) 
 &\aleq \sum_{k\in \N\cup\{0\}}\sum_{\ell\in\N\cup\{0\}}\int_{A(X,k,r)}\int_{A(X,\ell,r)} \abs{\Phi(\xi,t) - \Phi(\zeta,t)}\frac{r}{(2^k r)^{n+1}}\frac{r}{(2^\ell r)^{n+1}}\dif \xi \dif \zeta\\
 &\aleq \sum_{k\in \N\cup\{0\}}\sum_{\ell\in\N\cup\{0\}}2^{-k} 2^{-\ell}\barint_{A(X,k,r)}\barint_{A(X,\ell,r)} \abs{\Phi(\xi,t) - \Phi(\zeta,t)}\dif \xi \dif \zeta\\
 &\aleq \sum_{k\in \N\cup\{0\}}\sum_{\ell\in\N\cup\{0\}}2^{-k} 2^{-\ell} \brac{\barint_{A(X,k,r)}\barint_{A(X,\ell,r)} \abs{\Phi(\xi,t) - \Phi(\zeta,t)}^n\dif \xi \dif \zeta}^\frac{1}{n}\\
 &\aleq \sum_{k\in \N\cup\{0\}}\sum_{\ell\in\N\cup\{0\}}2^{-k} 2^{-\ell} \|\nabla \Phi(\cdot,t)\|_{L^n(\S^n)} \aleq \eps^{\frac 1n}.
 \end{split}
\]
Thus, $\left\| \dist(\Phi^h,\n) \right\|_{L^{\infty}(B^{n+1}\times[-1,1]^k)} \aleq \eps^{1/n}$.


Let $\pi_{\n}\colon B_{\delta}(\n) \to \n$ denote the nearest point projection. We then define
\[
 \Psi\coloneqq \pi_{\n} \circ \Phi^h \colon \B^{n+1} \times [-1,1]^k \to \n.
\]
By construction, $\Psi$ is continuous and satisfies
$
\Psi  \rvert_{\S^n\times [-1,1]^k} = \Phi.
$
Hence, we have $\left[\Psi \rvert_{\S^n\times [-1,1]^k} \right] = [\Phi]$ in $\pi_{k+n}(\n)$. On the other hand, the map $\Psi\rvert_{\S^n\times [-1,1]^k}$ is contractible, therefore $\left[\Psi\rvert_{\S^n\times [-1,1]^k} \right]=[\Phi]=0$. This is a contradiction with the assumption that $\alpha=[\Phi]$ is nontrivial.
\end{proof}

\end{document}
